%============================================================================%
%  INFOGRAPHICS CV -- Claudio Chagas
%  Based on Jan Kuster's Infographics CV template (MIT Licence)
%  Self-contained: g/ helpers inlined, icons via fontawesome (no PNG files)
%============================================================================%

\documentclass[10pt]{article}

%--- ENCODING ----------------------------------------------------------------
\usepackage[utf8]{inputenc}

%--- LOGIC -------------------------------------------------------------------
\usepackage{xifthen}
\usepackage{calc}

%--- FONT --------------------------------------------------------------------
\usepackage[default]{raleway}
\renewcommand*\familydefault{\sfdefault}
\usepackage[T1]{fontenc}
\usepackage{moresize}

%--- ICONS (fontawesome package -- no PNG files needed) ----------------------
\usepackage{fontawesome5}
\usepackage{textcomp}
\renewcommand{\copyright}{\textcopyright}

%--- PAGE LAYOUT -------------------------------------------------------------
\usepackage{fancyhdr}
\fancyhf{}
\renewcommand{\headrule}{}
\lfoot{\footnotesize Copyright \copyright\ Claudio Chagas, \today}

\usepackage[a4paper]{geometry}
\geometry{top=1cm, bottom=1cm, left=1cm, right=1cm}
\pagestyle{fancy}
\setlength{\headheight}{0.0pt}
\lhead{}\rhead{}\chead{}
\setlength{\parindent}{0mm}

%--- TABLE -------------------------------------------------------------------
\usepackage{array}
\newcolumntype{x}[1]{>{\raggedleft\hspace{0pt}}p{#1}}

%--- GRAPHICS ----------------------------------------------------------------
\usepackage{graphicx}
\usepackage{wrapfig}
\usepackage{tikz}
\usetikzlibrary{shapes, backgrounds}

\newcommand{\vcenteredinclude}[1]{\begingroup
\setbox0=\hbox{\includegraphics{#1}}%
\parbox{\wd0}{\box0}\endgroup}

\newcommand*{\vcenteredhbox}[1]{\begingroup
\setbox0=\hbox{#1}\parbox{\wd0}{\box0}\endgroup}

%--- ICON COMMANDS (fontawesome, replaces PNG system) ------------------------
\newcommand{\icon}[2]{\colorbox{thirdcol}{\makebox(#2,#2){\textcolor{sectcol}{\csname fa#1\endcsname}}}}
\newcommand{\icontext}[3]{%
	\vcenteredhbox{\icon{#1}{#2}} \vcenteredhbox{\textcolor{textcol}{#3}}%
}

%--- COLOURS -----------------------------------------------------------------
\usepackage{color}
\definecolor{maincol}{RGB}{255,150,0}
\definecolor{secondcol}{RGB}{0,178,255}
\definecolor{thirdcol}{RGB}{0,80,130}
\definecolor{fourthcol}{RGB}{0,100,160}
\definecolor{fifthcol}{RGB}{0,120,200}
\definecolor{sixthcol}{RGB}{0,80,130}
\definecolor{bgcol}{RGB}{190,220,255}
\definecolor{textcol}{RGB}{0,50,80}
\definecolor{sectcol}{RGB}{255,255,255}
\pagecolor{bgcol}

%--- PAGE STYLE --------------------------------------------------------------
\renewcommand{\headrulewidth}{0pt}
\renewcommand{\footrulewidth}{0pt}
\renewcommand{\thepage}{}
\renewcommand{\thesection}{}

%--- TIKZ ARROWS -------------------------------------------------------------
\newcommand{\tzlarrow}{(0,0) -- (0.2,0) -- (0.3,0.2) -- (0.2,0.4) -- (0,0.4) -- (0.1,0.2) -- cycle;}
\newcommand{\tzrarrow}{(0,0.2) -- (0.1,0) -- (0.3,0) -- (0.2,0.2) -- (0.3,0.4) -- (0.1,0.4) -- cycle;}

\newcommand{\larrow}[1]
{\begin{tikzpicture}[scale=0.58]
	\filldraw[fill=#1!100,draw=#1!100!black] \tzlarrow
\end{tikzpicture}}

\newcommand{\rarrow}[1]
{\begin{tikzpicture}[scale=0.58]
	\filldraw[fill=#1!100,draw=#1!100!black] \tzrarrow
\end{tikzpicture}}

%============================================================================%
%  INLINED CHART HELPERS (from g/ directory -- Jan Kuster MIT)
%============================================================================%

%--- PIE CHART ---------------------------------------------------------------
%counters for chart loop
\newcounter{a}
\newcounter{b}
\newcounter{c}

% draw a slice for a chart
% param 1: Circle form - 90 = quarter, 180 = half, 360 = full
% param 2: scale default=1 (scales only chart, not label text)
% param 3: border color
% param 4: label text color
% param 5: label bg color
% param 6:
\newenvironment{piechart}[5] {

	% draw a slice for a chart
	% param 1: value x of 100
	% param 2: label text
	% param 3: fill color
	% param 4:
	% param 5:
	% param 6:
	\newcommand{\slice}[3] {

		\setcounter{a}{\value{b}}
		\addtocounter{b}{##1}

		%set from angle point
		\pgfmathparse{\thea/100*#1}
	  	\let\pointa\pgfmathresult

		%set toanglepoint
		\pgfmathparse{\theb/100*#1}
	  	\let\pointb\pgfmathresult

		%set midangle
	 	\pgfmathparse{0.5*\pointa+0.5*\pointb}
	  	\let\midangle\pgfmathresult
		
		% draw the slice
	  	\filldraw[fill=##3!100,draw=#3!100, line width=2pt ] (0,0) -- (\pointa:#2) arc (\pointa:\pointb:#2) -- cycle;

	  	% draw label
	  	\node[label=\midangle:\colorbox{#5}{\textcolor{#4}{##2}}] at (\midangle:#2) {};

		\filldraw[fill=#3,draw=none] (0,0) circle (#2/2);
	}

	% execute commands
	\setcounter{a}{0}
	\setcounter{b}{0}
	\begin{tikzpicture}
}
{\end{tikzpicture}}

%--- BAR CHART ---------------------------------------------------------------
\newcounter{barcount}


% draw a bar chart
% param 1: width
% param 2: height
% param 3: border color
% param 4: label text color
% param 5: label bg color
% param 6: cat 1 color
\newenvironment{barchart}[8]{

	\newcommand{\barwidth}{0.35}
	\newcommand{\barsep}{0.2}

	% param 1: overall percent
	% param 2: label
	% param 3: cat 1 percent
	% param 4: cat 2 percent
	% param 5: cat 3 percent
	\newcommand{\baritem}[5]{

		\pgfmathparse{##3+##4+##5}
		 \let\perc\pgfmathresult

		\pgfmathparse{#2}
		 \let\barsize\pgfmathresult
	
		\pgfmathparse{\barsize*##3/100}
		 \let\barone\pgfmathresult
	
		\pgfmathparse{\barsize*##4/100}
		 \let\bartwo\pgfmathresult
	
		\pgfmathparse{\barsize*##5/100}
		 \let\barthree\pgfmathresult

		\pgfmathparse{(\barwidth*\thebarcount)+(\barsep*\thebarcount)}
		 \let\barx\pgfmathresult

		\filldraw[fill=#6, draw=none] (0,-\barx) rectangle (\barone,-\barx-\barwidth);
		\filldraw[fill=#7, draw=none] (\barone, -\barx) rectangle (\barone+\bartwo,-\barx-\barwidth);
		\filldraw[fill=#8, draw=none] (\barone+\bartwo,-\barx ) rectangle (\barone+\bartwo+\barthree,-\barx-\barwidth);

		\node [label=180:\colorbox{#5}{\textcolor{#4}{##2}}] at (0,-\barx-0.175) {};
		\addtocounter{barcount}{1}
	}
	\begin{tikzpicture}
	\setcounter{barcount}{0}
	
}
{\end{tikzpicture}}

%--- BUBBLE CHART ------------------------------------------------------------
\newcommand{\bubble}[5]{
	\definecolor{tmpcol}{RGB}{50,50,#5}
	% slice
  	\filldraw[fill=thirdcol,draw=none] (#1,0.5) circle (#3);

  	% outer label
  	\node[label=\textcolor{textcol}{#4}] at (#1,0.7) {};
}

\newcommand{\bubbles}[2]{
	%reset counters
	\setcounter{a}{0}
	\setcounter{c}{150}
	\begin{tikzpicture}[scale=3]
	\foreach \p/\t in {#1} {
	    	\addtocounter{a}{1}
	    	\bubble{\thea/2}{\theb}{\p/25}{\t}{1\p0}
	}
	\end{tikzpicture}
}

%--- SQUARE CHART ------------------------------------------------------------
\newcommand{\squares}[2]{
	%reset counters
	\setcounter{a}{0}
	\setcounter{b}{0}
	\setcounter{c}{50}
	\begin{tikzpicture}[scale=3]
	\foreach \p/\t in {#1} {
		\setcounter{a}{\value{b}}
	    	\addtocounter{b}{\p}
	    	\square{\thea/100*#2}{\theb/100*#2}{\p\%}{\t}{thirdcol}
	}
	\end{tikzpicture}
}

\newcommand{\square}[5] {
 	\pgfmathparse{#1+0.5*(#2-#1)}
  	\let\midangle\pgfmathresult

	\draw[draw=sectcol] (0.4, \midangle) -- (0.6,\midangle);
	% slice
  	\filldraw[fill=#5!100,draw=bgcol!100, line width=3pt] (0,#1) -- (0.5,#1) -- (0.5,#2) -- (0,#2) -- cycle;
  	% outer label
  	\node[label=360:\colorbox{sectcol}{\textcolor{textcol}{#4}}] at (0.55,\midangle) {};
}

%--- TIMELINE ----------------------------------------------------------------
% define global counters
\newcounter{yearcount}


\newcounter{leftcount}

% env cvtimeline
%
% creates a vertical cv timeline
%
% param 1: start year
% param 2: end year
% param 3: overall width
%param 4: overall height
\newenvironment{cvtimeline}[4]{
	
	\newcommand{\cvcategory}[2]{
		\node[label=\mbox{\colorbox{##1}{\strut\hspace{2pt}}\colorbox{white}{\textcolor{textcol}{##2}}}] at (0,-5) {}; %start year
	}

	\newcommand{\bxwidth}{4.5}
	\newcommand{\bxheight}{2}


	% creates a stretched box as cv entry headline followed by two paragraphs about 
	% the work you did
	% param 1:	event start month/year
	% param 2:	event end month/year
	% param 3:  event name
	% param 4:	institution (where did you work / study)
	% param 5:	what was your position
	% param 6:	color
	% param 7: level (position, use minus for left placement)
	\newcommand{\cvevent}[7] {
		
		\foreach \monthf/\yearf in {##1} {
			\foreach \montht/\yeart in {##2} {

				\pgfmathparse{#3/\fullrange*((\yearf-#1)+(\monthf/12))}
  				\let\startexp\pgfmathresult
				\pgfmathparse{#3/\fullrange*((\yeart-#1)+(\montht/12))}
  				\let\endexp\pgfmathresult
				\pgfmathparse{1/(\endexp-\startexp+1)}
  				\let\lenexp\pgfmathresult
				\pgfmathparse{0.5*\endexp+0.5*\startexp}
  				\let\midexp\pgfmathresult

				%\filldraw[fill=##6,draw=none, opacity=0.9] (-0.15-##7,\startexp) rectangle (-0.15-##7-0.5,\endexp);
				\draw[draw=##6, line width=1.5pt] (0, \startexp) -- (1,\startexp);

				\node[label={[label distance=0]0:\colorbox{##6}{\strut}\colorbox{white}{\textcolor{gray}{##1}\hspace{3pt}\textcolor{textcol}{##3}}}] at (0.5, \startexp) {};
			}
			\addtocounter{leftcount}{1}
		}
	}

	%--------------------------------------------------------------------------------------
	%	BEGIN
	%--------------------------------------------------------------------------------------

	\begin{tikzpicture}

	\setcounter{leftcount}{1}

	%calc fullrange= number of years
 	\pgfmathparse{(#2-#1)}
  	\let\fullrange\pgfmathresult
	\draw[draw=textcol,line width=4pt] (0,0) -- (0,#3) ;	%the timeline

	%for each year put a horizontal line in place
	\setcounter{yearcount}{1}
	\whiledo{\value{yearcount} < \fullrange}{
		\draw[fill=white,draw=textcol, line width=2pt]  (0,#3/\fullrange*\value{yearcount}) circle (0.1);
		\stepcounter{yearcount}
	}

	%start year
	\filldraw[fill=white!100,draw=textcol,line width=3pt] (0,-0.5) circle (0.5);
	\node[label=\textcolor{textcol}{\textbf{\small#1}}] at (0,-0.85) {}; 

	%end year
	\filldraw[fill=white!100,draw=textcol,line width=5pt] (0,#3+0.75) circle (0.75);
	\node[label=\textcolor{textcol}{\textbf{\large#2}}] at (0,#3+0.42) {}; 


	
}%end begin part of newenv
{\end{tikzpicture}}

%--- FACT BUBBLE -------------------------------------------------------------
% draw a circle with facts
% param 1: fact text
% param 2: scale default=1 (scales only chart, not label text)
% param 3: big border color
% param 4: second border color
% param 5: label bg color
\newcommand{\factbubble}[5]{
	\begin{tikzpicture}
	\pgfmathparse{#2*2}
	\let\pbxwidth\pgfmathresult
		\filldraw[fill=#3,draw=none] (0,0) circle (#2 * 1.5);
		\filldraw[fill=#5,draw=#4, line width=3.5pt] (0,0) circle (#2 * 1.2);
		\node at (0,0) {
			\parbox{\pbxwidth cm}{
				\begin{center}	
					#1
				\end{center}
			}
		};
	\end{tikzpicture}
}


%============================================================================%
%  CUSTOM SECTION COMMANDS
%============================================================================%

\newcommand{\cvsection}[2]{%
	\textcolor{sectcol}{\uppercase{\textbf{#1}}}%
}

\newcommand{\cvsect}[4]{%
	\textcolor{#3}{\hrule}%
	\colorbox{#3}{{\cvsection{#1}{#4}}}%
}

\newcommand{\metasection}[2]{%
	\begin{tabular*}{1\textwidth}{ l l }
		#1&#2\\[12pt]
	\end{tabular*}%
}

\newcommand{\cveventmeta}[2]{%
	\mbox{\mystrut \hspace{87pt}\textit{#1}}\\
	#2%
}

\newcommand{\mystrut}{\rule[-.3\baselineskip]{0pt}{\baselineskip}}
\newcommand{\legend}[3]{\parbox{#2}{\textcolor{#3}{\rule{#2}{4pt}}\\#1}}
\newcommand{\lorem}{Lorem ipsum dolor sit amet.}

%============================================================================%
%  DOCUMENT
%============================================================================%
\begin{document}

\pagestyle{fancy}

%--- TITLE HEADLINE ----------------------------------------------------------
\mystrut
\vspace{-12pt}

\begin{tabular*}{1\textwidth}{c c c}
	\parbox[c]{0.4\linewidth}{
		\colorbox{thirdcol}{\huge{}{\textcolor{white}{\textbf{\uppercase{Claudio Chagas}}} }}\\
		\Large{\textcolor{thirdcol}{\textsc{Senior Localization Specialist | EN$\rightarrow$PT-BR}}}\\
	}&
	\parbox{0.25\textwidth}{
	\icontext{MapMarker}{10pt}{Rio de Janeiro, Brazil}\\
	\icontext{Mobile}{10pt}{+55 21 999 5602 54}\\
	\icontext{Envelope}{10pt}{aviatorone@gmail.com}\\
	}&
	\parbox{0.3\textwidth}{
	\icontext{MousePointer}{10pt}{www.logostranslate.wordpress.com}\\
	\icontext{Github}{10pt}{github.com/claudchagas}\\
	\icontext{Linkedin}{10pt}{linkedin.com/in/claudiochagas}\\
	}
\end{tabular*}

\small
\vspace{16pt}

\begin{minipage}{0.59\textwidth}

%--- FACTS -------------------------------------------------------------------
	\mbox{
		\parbox[c][3cm][c]{0.29\textwidth}{
		\textcolor{textcol}{My journey in translation began with volunteering for global non-profits, grounded in a Bachelor's degree in Publishing from Oxford. I now serve enterprise clients including OpenAI, Mailchimp, Keeper Security, Lucid Software, SAP, and Equinix through Smartling and Mojito.}
		}
		\hspace{10pt}
		\parbox[c][3cm][c]{0.32\textwidth}{
			\begin{center}
			\factbubble{\huge{\textcolor{sectcol}{\textbf{B.A.}}}\\\small{\textcolor{sectcol}{\textbf{Publishing}}}}{1}{maincol}{sectcol}{thirdcol}\\
			\textcolor{textcol}{latest degree from}\\
			\textcolor{textcol}{\textbf{Oxford Brookes University}}
			\end{center}
		}
		\hspace{10pt}
		\parbox[c][3cm][c]{0.29\textwidth}{
		\textcolor{textcol}{I specialise in technical UI localisation, AI platform content, SaaS documentation, and marketing transcreation for global technology enterprises targeting Brazilian audiences.}
		}
	}
	\vspace{34pt}

%--- SKILLS AND TECHNOLOGIES -------------------------------------------------
	\cvsect{Skills and Technologies}{0.49}{thirdcol}{textcol}\\[4pt]
	\mbox{
		\parbox[b][150pt][c]{0.35\textwidth}{
			\textcolor{textcol}{I combine deep linguistic expertise with proprietary workflow infrastructure, including Python tooling for glossary automation and Mojito skill files for enterprise accounts.}\\

			\cvsect{Languages}{0.49}{thirdcol}{textcol}\\[4pt]
			\icontext{Language}{14pt}{\colorbox{maincol}{Portuguese (native)}}\\
			\icontext{Language}{14pt}{\colorbox{maincol}{English (C2)}}\\
		}

		\begin{piechart}{360}{2}{bgcol}{textcol}{sectcol}
			\slice{16}{Copyediting}{thirdcol}
			\slice{6}{Subtitling}{fourthcol}
			\slice{10}{Personal Admin}{fifthcol}
			\slice{16}{Proofreading}{secondcol}
			\slice{54}{Translation}{maincol}
		\end{piechart}\\
	}
	\begin{center}
	\begin{tikzpicture}
		\draw[draw=sectcol,dashed, opacity=0.5] (4,0) -- (-4,0);
	\end{tikzpicture}
	\end{center}

	\mbox{\hspace{-14pt}
		\begin{barchart}{10}{5.5}{sectcol}{textcol}{sectcol}{maincol}{secondcol}{thirdcol}
			\baritem{90}{Smartling / Mojito}{100}{0}{0}
			\baritem{80}{SDL Trados Studio}{0}{0}{70}
			\baritem{75}{Python / GitHub}{0}{65}{0}
			\baritem{65}{Google Suite}{0}{0}{60}
			\baritem{60}{LaTeX}{0}{0}{50}
		\end{barchart}
		\hspace{10pt}
		\parbox[b][100pt][c]{0.3\textwidth}{\textcolor{textcol}{I work across Smartling and Mojito for enterprise clients including OpenAI, Keeper Security, Lucid Software, and SAP, delivering UI strings, technical documentation, and marketing content.}}
	}
	\begin{center}
	\begin{tikzpicture}
		\draw[draw=sectcol,dashed, opacity=0.5] (4,0) -- (-4,0);
	\end{tikzpicture}
	\end{center}
	\begin{center}
	\mbox{
		\parbox[b][2cm][c]{3cm}{
			\textcolor{textcol}{I work with a specialised toolset built around enterprise localisation platforms and custom automation.}
		}\hspace{12pt}
		\bubbles{6/Smartling, 5/Mojito, 4/Trados, 4/Python, 3/LaTeX}{\cvsection{Technologies}}
	}
	\end{center}

%--- ACTIVITIES --------------------------------------------------------------
	\cvsect{Activities}{0.49}{thirdcol}{textcol}\\[20pt]
	\mbox{
		\parbox[b][3cm][c]{3cm}{
			\factbubble{\HUGE{\textcolor{sectcol}{\textbf{15+}}}\\\large{\textcolor{sectcol}{\textbf{years}}}}{0.85}{maincol}{sectcol}{thirdcol}\\
			\textcolor{textcol}{of experience in the}\\
			\textcolor{textcol}{\textbf{\mbox{Translation Industry}}}
		}

		\parbox[b][3cm][c]{4cm}{
			\textcolor{textcol}{I enjoy cycling to help maintain a healthy work-life balance and escape the city's hustle and bustle.}
		}

		\squares{20/Meditation,40/Biking,10/News,20/Music}{1}
	}
\end{minipage}
\begin{minipage}{0.05\textwidth}
	\begin{center}
		\begin{tikzpicture}
			\draw[draw=sectcol,dashed, opacity=0.5] (0,-12) -- (0,12);
		\end{tikzpicture}
	\end{center}
\end{minipage}
\begin{minipage}{0.4\textwidth}

%--- EXPERIENCE AND EDUCATION ------------------------------------------------
\cvsect{Experience and Education}{0.4}{thirdcol}{textcol}\\[16pt]

\hspace{60pt}\mbox{\legend{Experience}{1.8cm}{maincol} \legend{Events}{1.1cm}{secondcol} \legend{Education}{1.5cm}{thirdcol}}
\vspace{-40pt}
\begin{center}

\begin{cvtimeline}{2000}{2026}{21}{\linewidth}
\cvevent{8/2000}{7/2003}{Bachelor Studies, UK}{Oxford Brookes}{Master Thesis: Why Mass-Market Paperback Publishing is Still to Catch on in Brazil}{thirdcol}{0}
\cvevent{1/2001}{12/2002}{Translator/Copy-editor}{Viva Network}{Worked on the Portuguese edition of a facilitator's handbook and four other language versions of a children worker's pack}{maincol}{1}
\cvevent{6/2001}{7/2001}{Student Placement}{Oxford Brookes}{Worked as the administrator of the British Book Design and Production Awards \& BPIF for two consecutive years, in 2001 and 2002}{maincol}{1}
\cvevent{6/2002}{7/2002}{Student Placement}{Oxford Brookes}{Worked as the administrator of the British Book Design and Production Awards \& BPIF for two consecutive years, in 2001 and 2002}{maincol}{1}
\cvevent{3/2003}{3/2003}{The London Book Fair}{LBF}{The London Book Fair is the global marketplace for rights negotiation and the sale and distribution of content across print, audio, TV, film and digital channels}{secondcol}{0}
\cvevent{10/2003}{10/2003}{Ashmolean Museum Oxford}{University of Oxford}{Worked as Part-time Sales and Display Assistant at the Ashmolean Museum Oxford book and gift shop}{maincol}{1}
\cvevent{6/2004}{6/2004}{Rio Moto-GP 2004}{University of Oxford}{PA to the National Press Officer}{secondcol}{0}
\cvevent{4/2005}{5/2005}{Lexus Brochures Translation}{Just Customer Communication}{Proofreading and editing translated copy and managing electronic copy-editing}{maincol}{1}
\cvevent{9/2007}{11/2008}{Real Estate News Translation}{Dream Homes World Wide}{Work for an international real estate developer providing English translations from Portuguese and Spanish documents}{maincol}{1}
\cvevent{8/2008}{9/2008}{User Interface Translation}{G.host Inc.}{Translation of the G.ho.st Virtual Computer user interface strings (Client Module -- 14,500+ words) into Brazilian Portuguese}{maincol}{1}
\cvevent{10/2009}{10/2009}{Blog Articles Translation}{Sivers.org}{Translation of blog articles into Brazilian Portuguese for Derek Sivers, founder of CD Baby}{maincol}{1}
\cvevent{3/2010}{3/2010}{Business Article Translation}{London Business School}{Translation of an article from the Brazilian business magazine Epoca into English, on the legacy of Jorge Paulo Lemann, the founder of Bank Garantia}{maincol}{1}
\cvevent{7/2010}{9/2010}{Admin Interface Translation}{SafeKidZone}{String translation for backend database of SafeKidZone and SafeTREC website}{maincol}{1}
\cvevent{3/2011}{4/2011}{Project Research Assistant}{University of Bremen}{Conduct Internet-based research on the presence of automakers on the social media landscape in Brazil, and translate relevant content into English}{maincol}{1}
\cvevent{9/2011}{9/2011}{Website Translation}{Fortune3, Inc.}{Translate 700 strings (approx. 5000 words) of an ecommerce website into Portuguese}{maincol}{1}
\cvevent{1/2009}{12/2012}{Communications Assistant}{Alpha International}{Assist the Relationship Manager and other managers of the national and international teams on a range of projects, from Jan 2009 to Dec 2012}{maincol}{1}
\cvevent{11/2012}{11/2012}{Alpha Conference Aparecida}{Alpha International}{Help set up the venue, providing an essential service for the conference and seminar sessions}{secondcol}{0}
\cvevent{7/2013}{7/2013}{World Youth Day 2013}{Alpha International}{Help find local suppliers, providing essential services for the event}{secondcol}{0}
\cvevent{1/2014}{1/2014}{Executive Assistant}{Heathfield}{Assist the Director with a range of roles to support business operations}{maincol}{1}
\cvevent{8/2016}{8/2016}{Rio 2016 Summer Olympics}{Organising Committee for the Rio 2016}{Help run the Venue Media Centre, providing an essential service and helping the press}{secondcol}{0}
\cvevent{6/2017}{10/2017}{Subtitles Translation}{ZOO Digital Group plc}{Provide subtitle translation services for films or any other localization-related tasks}{maincol}{1}
\cvevent{5/2019}{6/2024}{Senior Translator \& Localization Specialist}{Smartling / Mojito}{Delivering EN$\rightarrow$PT-BR localisation for enterprise clients including OpenAI, Mailchimp, Keeper Security, Lucid Software, SAP, and Equinix --- spanning AI platform UI, SaaS documentation, and marketing content}{maincol}{1}
\end{cvtimeline}
\end{center}
\end{minipage}

\end{document}